\section{Descomposición cíclica y cocientes de determinantes}
Los proyectores cíclicos
\[
P_3=\tfrac14(I+S^3+S^6+S^9),\qquad
P_4=\tfrac13(I+S^4+S^8)
\]
tienen imágenes de dimensiones tres y cuatro. Si $E=P_3-P_4$, entonces
\begin{equation}
\frac{\Tr E}{12}=-\frac1{12}.
\label{eq:minus12}
\end{equation}
El escalar procede de una diferencia de rangos normalizada, no de una
suma ordinaria de enteros positivos.

\begin{theorem}[Representación del lector de canal mediante un cociente de determinantes]
Sean $\mathcal H_3=\operatorname{Fix}(S^3)$ y
$\mathcal H_4=\operatorname{Fix}(S^4)$. Para $0<s<1$,
\begin{equation}
F(s)=\frac{\det_{\mathcal H_3}(I-sS)}{\det_{\mathcal H_4}(I-sS)}
=\frac{1-s^3}{1-s^4}
=\frac{1+s+s^2}{1+s+s^2+s^3}.
\label{eq:det}
\end{equation}
Su prolongación continua en $s=1$ vale $3/4$.
\end{theorem}
\begin{proof}
La restricción de $S$ a cada imagen es un ciclo de tres o cuatro
posiciones, respectivamente. Sus polinomios característicos dan
$\det(I-sS)=1-s^r$ para $r=3,4$. Al cancelar el factor común $1-s$,
queda la última expresión, regular en $s=1$.
\end{proof}

La sustitución $s=q^{30}$ da
$F(q^{30})=(1-q^{90})/(1-q^{120})$.
El mismo par de espacios que produce \eqref{eq:minus12}
produce el cociente de canal. No se identifica por esta igualdad
una matriz de doce coordenadas con todo el estado TPK.

\subsection{Información orientada y sector complejo}
En el sector antipodal $V_-=\operatorname{Ran}((I-S^6)/2)$,
el operador $J=S^3$ satisface $J^2=-I$ y $J^{\mathsf T}=-J$.
Con $Q=P_4(I-S^6)/2$ se tiene
\[
\operatorname{rank}Q=2,\qquad E|_{V_-}=-Q,\qquad B|_{V_-}=5I-3Q.
\]
En $\operatorname{Ran}Q$, $S^2=-I$ y el determinante es $1+s^2$.
Así,
\[
F(s)=\frac{1+s+s^2}{1+s}\,\frac1{1+s^2}.
\]
La inversión de orientación cambia el signo de la forma
$\omega(u,v)=\langle Ju,v\rangle$, pero conserva $B$ y $F$.
La evaluación escalar no distingue todos los transportes orientados.

La traza transportada produce además
\[
a_n^{\rm tr}=\Tr(S^nE)=3\mathbf1_{3\mid n}-4\mathbf1_{4\mid n},
\]
y, para $|s|<1$,
\[
-\log F(s)=\sum_{n\ge1}\frac{a_n^{\rm tr}}n\,s^n.
\]
Esta sucesión de momentos enlaza la descomposición cíclica con la
representación espectral del lector. La naturaleza aritmética de valores
particulares, como la constante de Catalán, requiere además el análisis
aritmético correspondiente.
